\documentclass[11pt]{article}
\usepackage[margin=1in]{geometry}
\usepackage{amsmath,amssymb,amsthm,hyperref}
\newtheorem{theorem}{Theorem}\newtheorem{lemma}[theorem]{Lemma}
\newtheorem{corollary}[theorem]{Corollary}
\title{Analytic interpolation of the physical moment pair on closed disks}
\author{Internal research note}\date{18 November 2026}
\begin{document}\setlength{\emergencystretch}{1em}\maketitle
Fix a prime $p$ and put $q=p-1$ and $w=(p-1)!$.
Use the polynomials $W_n,H_n$ of
\href{PHYSICAL_PRIME_MOMENT_PAIR_RECURRENCES.tex}
{the primitive moment pair}, \texttt{lem:pair}.
In particular $H_n=W_{n+1}+qW_n$ and $W_n(0)=1$.
For $n\ge2$, their \texttt{cor:critical} identifies a common
root of $W_n,W_{n+1}$ with a multiple root of $H_n$.

Let $F/\mathbb Q_p$ be finite, let $y\in\mathcal O_F^\times$ and choose
$0\ne\alpha\in\mathfrak m_F$. We study the fixed closed disk
\[
 \mathbb D=\{y+\alpha X:X\in\mathbb C_p,\ |X|_p\le1\}.
\]
The disk lies inside the residue class of $y$.
Enlarge $F$ unramifiedly to a field $L$ whose residue field
$\mathbb F_{p^f}$ splits
$X^p+\bar y X^{p-1}-1$. Set $m=p^f-1$ and use Gauss norms
on the Tate algebras over $L$.


\begin{theorem}[Uniform residue multiplicities]\label{thm:residue-multiplicity}
For every $n\ge1$, a unit root of $\bar W_n$ has multiplicity at
most $p-1$, and this bound is attained at every unit residue
for some terminal indices. A unit root of $\bar H_n$ has
multiplicity divisible by $p$. At a common unit root of
$\bar H_n,\bar W_n$, the multiplicity in $\bar W_n$ is at
most $p-2$.

Let $y_0$ be an unramified integral lift of any nonzero residue
in $\overline{\mathbb F}_p$. Every root of $W_n$ in its residue
class has $v_p(Y-y_0)\ge1/(p-1)$ and tame ramification over
$\mathbb Q_p$. For a common root of the primitive moment pair
and $p>2$, the bound improves to $1/(p-2)$ and the ramification
index is at most $p-2$. When $p=2$ the primitive moment pair
has no common root.
\end{theorem}
\begin{proof}
All the residue identities first take place over $\mathbb F_p(Y)$.
The order-$p$ residue denominator for $W_n$ is
$1+Yz-z^p$, with numerator $(1-z)^{p-1}$ by
\texttt{prop:rational} at $K=1$.
At a simple pole $z=1/\lambda$ the coefficient of
$(1-\lambda z)^{-1}$ is
\[
 -\frac{\lambda(1-1/\lambda)^{p-1}}{Y}
                     =(1-1/\lambda)^{-1},
\]
since $(1-1/\lambda)^p=-Y/\lambda$.
Thus $\bar W_n=\sum_\lambda\lambda^n/(\lambda-1)$.
For $\bar H_n$ the rational trace formula gives
$\bar H_n=\sum_\lambda\lambda^n$. This last residue trace follows
at every prime from \texttt{lem:pair} and the $K=1$ denominator:
the generating series is $p-zD'(z)/D(z)$ in characteristic $p$,
with $D=1+Yz-z^p$.

Use a new variable $S$ and put $B=S^p-S$, $Y=B^{-(p-1)}$.
The reciprocal roots $\tau_i=1/\lambda_i$, $i\in\mathbb F_p$,
satisfy $\tau_i^p-Y\tau_i=1$ and are given by
\[
                  \tau_i=(S+i)^p/B .
\]
Indeed $(S+i)^{p^2}-(S+i)^p=B^p$, and the displayed
roots are distinct. Also $\lambda_i-1=-(S+i)^{1-p}$.
Consequently
\[
 \bar W_n(Y)=-B^n\sum_{i\in\mathbb F_p}(S+i)^{-pn+p-1},
 \qquad
 \bar H_n(Y)=B^n\left(\sum_{i\in\mathbb F_p}(S+i)^{-n}\right)^p.
 \tag{5}
\]
At any $S_0$ above a unit $Y$, $B(S_0)\ne0$ and
$Y'(S_0)=(p-1)B(S_0)^{-p}\ne0$, so multiplicities
agree in the two coordinates.
Write $k=-pn+p-1$. The first $p$ Taylor coefficients of
the sum for $\bar W_n$ at $S_0$ are
\[
 \binom{k}{j}\sum_i (S_0+i)^k (S_0+i)^{-j},
                  \qquad 0\le j\le p-1 .
\]
All $\binom{k}{j}$ are nonzero modulo $p$, since $k\equiv p-1$.
The $p$ distinct numbers $(S_0+i)^{-1}$ form a Vandermonde
matrix and the weights $(S_0+i)^k$ are units.
The $p$ coefficients therefore cannot all vanish.
The second formula in (5) gives the assertion for $\bar H_n$.
Moreover the last Vandermonde sum, at $j=p-1$, is
$\sum_i(S_0+i)^{-pn}$. It vanishes if $\bar H_n$ vanishes.
Then at least one of the first $p-1$ Taylor coefficients
is nonzero, giving the stronger common-root bound.

For sharpness of the separate $\bar W_n$ bound, take $S_0$
in $\mathbb F_{p^f}$ and choose $n$ with
$pn\equiv1\pmod{p^f-1}$. The first $p-1$ Vandermonde sums
are then sums of $(S_0+i)^j$ with $0\le j<p-1$.
Each vanishes by the sums of powers in $\mathbb F_p$.
The last sum is nonzero by the Vandermonde determinant.

If the residue root of $\bar W_n$ has multiplicity $e$,
then the coefficients of $W_n(y_0+Z)$ of degree below $e$
are divisible by $p$ and its degree-$e$ coefficient is a unit.
For $0<v_p(Z)<1/e$, the degree-$e$ term has strictly smaller
valuation than every other term, so there is no root.
Hensel factorization after monic normalization of $W_n$
gives the factor in this residue class of degree $e$.
Each root generates an extension of the unramified coefficient
field of degree at most $e\le p-1$, hence with ramification
index prime to $p$. At a common root, apply the same argument
with $e\le p-2$. For $p=2$, the common residue bound excludes
every unit residue, and \texttt{cor:critical} places every
possible common root at a unit.
\end{proof}


\begin{theorem}[A uniform second-precision trace]\label{thm:second-trace}
Let $p$ be odd and $S_0\in\mathbb F_{p^f}\setminus\mathbb F_p$.
Choose a lift $S$ in the integers of the unramified extension of
degree $f$, and put $B=S^p-S$ and $y_0=B^{-(p-1)}$.
Write $\omega_i$ for the Teichm\"uller lift of $S_0+i$.
Then for every $n\ge1$,
\[
 H_n(y_0)\equiv (-wB)^n\sum_{i\in\mathbb F_p}\omega_i^{-pn}
                                                 \pmod{p^2}.
 \tag{6}
\]
For every $1\le\ell\le p-2$ and positive integer $n$ with
$-pn\equiv\ell\pmod{p^f-1}$, the residue of $y_0$ is a
root of $\bar H_n,\bar W_n$ of exact multiplicity $\ell$ in
$\bar W_n$. Nevertheless the entire residue class contains
no common root of the actual primitive pair.
In particular the common residue multiplicity bound is sharp
at each unit residue.
\end{theorem}
\begin{proof}
The congruence of $H_n$ modulo $p^2$ with the trace of the
reciprocal polynomial $D_0(z)=(1-z)^p-wy_0z$ holds for all
odd primes by \texttt{prop:trace} of the
\href{PHYSICAL_PRIME_TOTAL_DEGREE_REDUCTION.tex}{physical prime-total-degree note}.
Its residue roots are the inverses of those used in (5).
Choose integer representatives $0,\ldots,p-1$ of $i$ and put
$x_i=S+i$. The element $x_i^p/(-wB)$ solves $D_0(z)=0$
modulo $p^2$. Indeed it reduces to $x_i^p/B$ and
$1-z\equiv-x_i/B\bmod p$ at that value.
Taking $p$th powers modulo $p^2$ gives
$(1-z)^p\equiv-x_i^p/B^p=w y_0z\bmod p^2$.
The roots of $D_0$ are simple modulo $p$, so these
representatives give all its roots modulo $p^2$ by Hensel's lemma.
Their inverses are $-wB/x_i^p$ modulo $p^2$. Since $x_i$ and
$\omega_i$ agree modulo $p$, their $p$th powers agree modulo
$p^2$, proving (6).

For the specified $\ell$, the $\bar H_n$ sum in (5) is a
sum of $\ell$th powers of $S_0+i$ at the center, and vanishes
since $\ell<p-1$. The Vandermonde exponent for $\bar W_n$
in (5) reduces to $\ell+p-1$ modulo $p^f-1$.
For Taylor index $j<\ell$, the sum uses
$(S_0+i)^{p+(\ell-1-j)}$.
Since $(S_0+i)^p=(S_0+i)+(S_0^p-S_0)$, its sum reduces
to sums of powers of exponent at most $\ell<p-1$, and is zero.
For $j=\ell$ the exponent is $p-1$ and the sum is $-1$.
The binomial coefficients of the Taylor expansion are units
for $j\le p-1$, giving the exact multiplicity.

To evaluate the lift, the polynomial identity
\[
        \sum_{i=0}^{p-1}(S+i)^{p\ell}
                    \equiv p(S^p-S)^\ell \pmod{p^2}
 \tag{7}
\]
holds for $1\le\ell\le p-2$.
Expand in powers of $i$. The exponent zero contributes
$pS^{p\ell}$. An exponent $pj$, $1\le j\le\ell$, has
$\sum_i i^{pj}\equiv0\bmod p^2$: each $i^p$ agrees
modulo $p^2$ with its Teichm\"uller lift, and the sum of
their $j$th powers is zero. Every exponent not divisible
by $p$ has binomial coefficient divisible by $p$.
Of these, only exponents $(p-1)j$ for $1\le j\le\ell$ have
a nonzero power sum modulo $p$, equal to $-1$.
For such an exponent,
\[
 \frac1p\binom{p\ell}{(p-1)j}
    =\frac{\ell}{(p-1)j}\binom{p\ell-1}{(p-1)j-1}
    \equiv (-1)^{j-1}\binom{\ell}{j}\pmod p
\]
by Lucas' formula. This proves (7).

Unramified Frobenius on $\sum_i\omega_i^\ell$ gives
$\sum_i\omega_i^{p\ell}\equiv pB^\ell\bmod p^2$ by (7).
As $B$ is a unit, the sum has valuation one.
Equation (6) now gives $v_p(H_n(y_0))=1$.
Every candidate common root is $Y=y_0+Z$ with
$v_p(Z)\ge1/(p-2)$ by Theorem~\ref{thm:residue-multiplicity}.
In $H_n(y_0+Z)$ the coefficients of $Z^k$ for $1\le k<p$
are divisible by $p$, because the residue multiplicity
is a positive multiple of $p$.
Every nonconstant term thus has valuation greater than one.
The constant term has valuation one and prevents cancellation.
Finally, each unit residue can be represented by choosing
a root $S_0$ of $(S^p-S)^{p-1}=1/\bar y$ in the algebraic
closure. The congruence classes of $n$ used above exist
since $p$ is coprime to $p^f-1$.
\end{proof}


\begin{corollary}[Exclusion of maximal common multiplicity]
\label{cor:maximal}
Let $p$ be odd. Every common unit residue where the multiplicity
in $\bar W_n$ is $p-2$ has an entire residue class free of
common roots of the actual primitive pair.
In particular the gap-two pair at $p=3$ has no common root.
For $p\ge5$ any actual common root has ramification index at
most $p-3$ and $v_p(Y-y_0)\ge1/(p-3)$ from its unramified lift $y_0$.
\end{corollary}
\begin{proof}
Retain the variable $S_0$ and Vandermonde weights
$A_i=(S_0+i)^{-pn+p-1}$ of (5).
A common residue of multiplicity $p-2$ satisfies the
$p-1$ independent homogeneous equations using rows
$j=0,\ldots,p-3,p-1$ of that Vandermonde matrix.
Its vector of weights is in the resulting one-dimensional kernel.
The vector with coordinates $(S_0+i)^{2p-3}$ is in the same
kernel by the Taylor sum calculation for $\ell=p-2$
in Theorem~\ref{thm:second-trace}.
Hence $(S_0+i)^{-pn}=c(S_0+i)^{p-2}$ for a fixed nonzero $c$.
Multiplicativity of the Teichm\"uller lifts gives
$\omega_i^{-pn}=\widetilde c\,\omega_i^{p-2}$ with a common
unit $\widetilde c$.
The proof of (7) shows that the sum in (6) has valuation one.
As before this excludes actual common roots in the entire
residue class. For $p=3$ the only positive common multiplicity
allowed by Theorem~\ref{thm:residue-multiplicity} is now excluded.
For $p\ge5$ the distance and degree bounds apply with $e\le p-3$.
\end{proof}



\begin{theorem}[A necessary point-value congruence]\label{thm:third-precision}
Let $p\ge5$ and let $y_0$ be any unramified integral lift of a
nonzero residue. If $H_n$ has a multiple root reducing to
$\bar y_0$, then
\[
                    v_p(H_n(y_0))\ge 3+v_p(n).
\]
In particular the order-$p$ trace modulo $p^3$ gives a
necessary test at every terminal index.
\end{theorem}
\begin{proof}
Write a hypothetical common root of $H_n,W_n$ in this class
as $Y=y_0+Z$. The case $Z=0$ gives $H_n(y_0)=0$.
Otherwise let $e$ be the multiplicity of this residue in
$\bar W_n$. It satisfies $1\le e\le p-3$ by
Corollary~\ref{cor:maximal}. The distance and degree bounds
give $v_p(Z)\ge1/e$ and place $v_p(Z)$ in the value group
of an extension of the unramified coefficient field with
ramification index at most $e$.

By the exact primitive pair identity,
\[
 \frac{d}{dX}\frac{H_n(y_0+X)}{(y_0+X)^n}
           =-pn\frac{W_n(y_0+X)}{(y_0+X)^{n+1}}.
\]
Write the quotient on the right without $-pn$ as
$\sum_{j\ge0}a_jX^j$. Its coefficients are integral and
$a_j$ is divisible by $p$ for $j<e$.
Termwise integration and $H_n(Y)=0$ give
\[
 \frac{H_n(y_0)}{y_0^n}
                =pn\sum_{k\ge1}\frac{a_{k-1}Z^k}{k}.
 \tag{8}
\]
The series converges for $v_p(Z)>0$, since $v_p(k)$
grows at most logarithmically in $k$.
The sum belongs to the unramified coefficient field by
the left side of (8).

For $k\le e$ the valuation of the $k$th summand is greater
than one by the divisibility of $a_{k-1}$.
For $k>e$ and $p\nmid k$ it is at least $k/e>1$.
For $p\mid k$ and $k\ge2p$, it is at least
$k/e-v_p(k)>1$. Indeed for $v_p(k)=1$ use $k/e\ge2p/e>2$,
and for $v_p(k)=j\ge2$ use $k/e>p^{j-1}\ge j+1$.
The only remaining index is $k=p$.
If its valuation were at most one, it would have to be
nonintegral: $v_p(a_{p-1})$ is integral and
$v_p(Z)\in(0,2/p]$, so $p v_p(Z)$ cannot be integral in a
value group with ramification index at most $e<p$.
That term would be uniquely minimal and the sum in (8)
would have nonintegral valuation, a contradiction.
Thus every term has valuation greater than one, and their
valuations tend to infinity. Taking valuations in (8)
gives $v_p(H_n(y_0))>2+v_p(n)$.
This is the claimed bound, since the point value is unramified.
The final trace test follows from \texttt{lem:cubic} of the
primitive moment pair note.
\end{proof}

\begin{theorem}[Analytic functions of the terminal index]\label{thm:interpolation}
There is an $s\ge0$, independent of the terminal order, such that
for each $r\bmod m$ and $0\le b<p^s$ there are
$F_0,F_1\in L\langle U,X\rangle$ with
\[
 F_i\!\left((n-b)/p^s,X\right)=W_{n+i}(y+\alpha X)
 \qquad (i=0,1)
\]
for all positive integers $n\equiv r\pmod m$, $n\equiv b\pmod{p^s}$.
The linear combination $\mathsf H=F_1+qF_0$ satisfies
\[
 (b+p^sU)\mathsf H-
 \frac{y+\alpha X}{\alpha}\partial_X\mathsf H
                   =p(b+p^sU)F_0 .
\]
\end{theorem}
\begin{proof}
Set $Y=y+\alpha X$ and $A^\circ=\mathcal O_L\langle X\rangle$.
The denominator $D_0=(1-z)^p-wYz$ has reciprocal
$(\lambda-1)^p-wY\lambda^{p-1}$. Its reduction at the disk
is $f_{\bar y}(\lambda)=\lambda^p+\bar y\lambda^{p-1}-1$.
Every residue root is a unit and simple since
$f_{\bar y}'=-\bar y\lambda^{p-2}$.
The reciprocal polynomial splits into $p$ analytic roots
$\lambda_i\in A^\circ$, all units with unit differences.
Indeed, start at the Teichm\"uller lift $\omega_i\in\mathcal O_L$ of
each residue root. Then the reciprocal polynomial has norm less
than one at $\omega_i$ and its derivative is a unit there.
Newton iteration in the complete Tate algebra gives a root with
$\|\lambda_i-\omega_i\|<1$. The differences remain units by
their distinct residues. Thus $D_0=\prod_i(1-\lambda_i z)$.

Use the elementary symmetric sums $E_j(k)$ from the
proof of \texttt{thm:linear} in the primitive moment pair note.
Set $a_p=1,2,3$ for $p=2,3,\ge5$, respectively. The proof gives
$\deg E_j\le pj$ and $E_j(\mathbb Z_p)\subseteq p^{a_pj}\mathbb Z_p$.
Write
\[
 Q=\frac{Yz}{(1-z)^p},\qquad
 G_j(z)=\frac{1}{1-z}\sum_{k\ge0}E_j(k)(wQ)^k.
\]
For each coefficient of $z$, $\sum_jG_j$ is finite and agrees
with the generating function of $W_n$ because $E_j(k)=0$ for $j>k$.

The binomial coefficients of $E_j(k)$ in its Newton expansion
are divisible by $p^{a_pj}$ by the finite difference formula.
Therefore $G_j$ is proper rational with denominator
$D_0^{pj+1}$ and numerator in $p^{a_pj}A^\circ[z]$.
Partial fractions give
\[
 G_j(z)=\sum_{i=1}^{p}\sum_{h=1}^{pj+1}
                           \frac{c_{j,i,h}}{(1-\lambda_i z)^h},
 \qquad c_{j,i,h}\in p^{a_pj}A^\circ .
 \tag{1}
\]
To see the bound, expand the numerator and the other factors
at $z=(1-T)/\lambda_i$. Each other reciprocal factor has a
unit constant term and integral Taylor coefficients there.
The principal part in $T$ gives the asserted coefficients.
The rational function is proper, so its principal parts give
the whole function.

Write $\lambda_i=\omega_i\kappa_i$ with $\|\kappa_i-1\|<1$.
Choose $s$ with all
$\|\kappa_i^{p^s}-1\|<p^{-1/(p-1)}$ and
\[
                  a_p-\frac{p}{p^s(p-1)}>0.
 \tag{2}
\]
Such an $s$ exists since a $p$th power increases any positive
lower valuation bound $\eta$ to at least $\min(1+\eta,p\eta)$.
For $N=b+p^sU$, interpolation of the power $\lambda_i^{n+i'-1}$
for $i'\in\{0,1\}$ is
\[
 \omega_i^{r+i'-1}\kappa_i^{b+i'-1}
             \exp\bigl(U\log(\kappa_i^{p^s})\bigr).
\]
The chosen bound gives convergence in the Tate algebra and
norm one. This agrees with the power for every integer $n$
in the given congruence class.

The coefficient of $z^{n+i'-1}$ in a term of (1) is its power
$\lambda_i^{n+i'-1}$ times
$\binom{n+i'+h-2}{h-1}c_{j,i,h}$.
For any integer $B$ and $t\ge0$, the Gauss norm bound
\[
 \left\|\binom{B+p^sU}{t}\right\|
 \le p^{\sum_{\ell>s}\lfloor t/p^\ell\rfloor}
 \le p^{t/(p^s(p-1))}
 \tag{3}
\]
holds. For each $\ell\le s$, at least $\lfloor t/p^\ell\rfloor$
of the $t$ consecutive numerator factors contribute a factor
of $p^\ell$ to the count of their valuations at level $\ell$.
Subtracting these counts from the valuation of $t!$ proves (3).
Consequently the whole $j$th contribution to each interpolant
has norm at most $p^{-j(a_p-p/(p^s(p-1)))}$.
Its sum over $j$ converges in $L\langle U,X\rangle$ and agrees
with the finite coefficient calculation at every prescribed $n$.
The differential identity holds at all such $n$ by the primitive
pair formula. The corresponding $U$ are dense in $\mathbb Z_p$.
An analytic function over $L$ vanishing at those $U$ vanishes
identically, coefficient by coefficient in $X$.
\end{proof}

\begin{corollary}[The same branches have explicit nonpositive values]
\label{cor:negative}
For an integer $u\ge0$ and the branch containing $n=-u$,
the continuation satisfies
\[
 F_0\bigl((-u-b)/p^s,X\bigr)=
 \sum_{0\le k\le u/q}(-1)^{pk}U_p(k)\binom{u+k}{pk}(y+\alpha X)^k .
 \tag{4}
\]
In particular $W_0$ in this continuation is one.
\end{corollary}
\begin{proof}
For a summand $(1-\lambda z)^{-h}$ in (1), the coefficient
of $z^{-u-1}$ in its expansion at infinity is
$(-1)^h\lambda^{-u-1}\binom{u}{h-1}$.
Its negative is exactly
$\lambda^{-u-1}\binom{-u+h-2}{h-1}$, the interpolated value.
At infinity $Q=O(z^{-q})$, so the coefficient of $z^{-u-1}$ in
the total series involves only $k\le u/q$. Since $E_j(k)$
vanishes for $j>k$, this also uses finitely many $j$.
Expansion of
$Y^k z^k(1-z)^{-pk-1}$ at infinity gives the negative of the
summand in (4), with weight $U_p(k)$.
\end{proof}

\begin{theorem}[Index sets for a fixed closed disk]\label{thm:local-set}
The set of integers $n\ge2$ for which $H_n$ has a multiple root
in $\mathbb D$ is the union of finitely many full arithmetic
progressions within $n\ge2$ and a finite set.
More precisely, on a finite refinement of the branches of
Theorem~\ref{thm:interpolation}, this set is the zero set of
a single analytic resultant in the branch coordinate. A full
progression occurs precisely when that resultant vanishes
identically on the branch.
\end{theorem}
\begin{proof}
Here are the required preparation and finiteness facts.
Over a Tate algebra in $U$, a series in $X$ of the form
$P_0+G$, where $P_0$ is monic of degree $d$ with lower coefficients
of norm at most one and $\|G\|<1$, is a unit times a monic
polynomial of degree $d$ with integral coefficients.
Division by $P_0$ on polynomials has quotient and remainder
of norm at most the input norm, by the monic division algorithm.
It extends by continuity to restricted series. To divide
$X^d$ by $P_0+G$, solve for the quotient by the iteration
$Q\mapsto\operatorname{quot}_{P_0}(X^d-GQ)$.
This is a contraction, and the fixed point differs from one
by norm less than one. The resulting remainder $R$ has
degree below $d$ and norm at most one. Thus $X^d-R=Q(P_0+G)$
with $Q$ a unit. In one variable this also proves the needed
zero finiteness: normalize a nonzero restricted series by its
last coefficient of maximal norm and apply the same argument.

Now fix an initial branch. Take $F_0$ from
Theorem~\ref{thm:interpolation}. It is nonzero since each
$W_n(0)=1$ and a nonzero polynomial remains nonzero after
the change of variable $Y=y+\alpha X$.
There are only finitely many $u\in\mathbb Z_p$ with
$F_0(u,X)$ identically zero, since any fixed nonzero
$X$-coefficient of $F_0$ has finitely many zeros.
At each such $u$ remove the maximal common power of $U-u$
from $F_0$. This division remains in the Tate algebra:
if $f(u)=0$ then $\sum_k a_k(U^k-u^k)/(U-u)$ converges
in norm when $f=\sum_k a_kU^k$, including with coefficients
in $L\langle X\rangle$. The common powers are bounded by the
multiplicities of the chosen nonzero coefficient. None of these
removed index points corresponds to an actual positive integer
on the branch. The quotiented function $F_0^*$ is nonzero as a
function of $X$ at every $u\in\mathbb Z_p$.

For each such $u$, normalize $F_0^*(u,X)$ by its last coefficient
of maximal norm, of index $d$. On a small enough residue
class of $u$ in $U$, the same constant scalar normalization
has all coefficients through degree $d$ of norm at most one,
the $d$th within norm less than one of one, and the tail
of norm less than one. This follows by substituting
$U=u+p^tV$, since the difference from evaluation at $u$
is bounded in Gauss norm by $p^{-t}\|F_0^*\|$ before
normalization. Compactness gives a finite refinement by
residue classes on which the preparation just proved applies.
Each refined class still corresponds to one arithmetic progression
of positive $n$.

On each refined class, let $P$ be the prepared monic polynomial
and let $R_1$ be the remainder of $F_1$ modulo $P$.
The determinant of multiplication by $R_1$ on the free module
with basis $1,X,\ldots,X^{d-1}$ modulo $P$ is analytic in the
refined index variable. For $d=0$ use determinant one.
After specialization, this determinant vanishes exactly when
$P$ and $R_1$ have a common root. All roots of $P$ are
in $|X|_p\le1$ since its coefficients are integral and it is monic.
The units and the removed index factors preserve common roots
at every actual positive integer on the branch. The determinant
thus has the asserted common-root interpretation. It is identically
zero or has finitely many zeros by one-variable preparation.
Finally use \texttt{cor:critical} of the primitive moment pair.
\end{proof}


\begin{corollary}[Index sets in a full unit residue class]\label{cor:full-residue}
Fix any nonzero $\bar y\in\overline{\mathbb F}_p$.
The set of positive terminal indices $n\ge2$ with a multiple
root of $H_n$ reducing to $\bar y$ is a finite union of full
arithmetic progressions within $n\ge2$ and a finite set.
\end{corollary}
\begin{proof}
For $p=2,3$ the set is empty by Theorem~\ref{thm:residue-multiplicity}
and Corollary~\ref{cor:maximal}.
For $p\ge5$ choose an unramified lift $y_0$ of $\bar y$ and extend
the coefficient field to contain $\alpha$ with $\alpha^{p-3}=p$.
By \texttt{cor:critical} of the primitive moment pair, a multiple
root of $H_n$ is a root of $W_n$.
Corollary~\ref{cor:maximal} confines every such root in the
full residue class to $|Y-y_0|\le|\alpha|$.
Theorem~\ref{thm:local-set} applies to this fixed closed disk,
which lies inside the residue class.
\end{proof}

A single index without a multiple root on a prepared branch rules
out the identically zero resultant there. The progressions and finite exceptions must still be excluded
for all unit residues to prove all-index squarefreeness.
\end{document}
